\documentclass[11pt]{article}
\usepackage[utf8]{inputenc}
\usepackage[T1]{fontenc}
\usepackage[margin=1in]{geometry}
\usepackage{amsmath,amssymb,mathtools,bm}
\usepackage{booktabs,array}
\usepackage{hyperref}
\usepackage{xcolor}
\usepackage{listings}

\hypersetup{colorlinks=true,linkcolor=blue,urlcolor=blue,pdftitle={EKF Explanation}}
\lstset{basicstyle=\ttfamily\small,columns=fullflexible,keepspaces=true,frame=single,breaklines=true}

\newcommand{\vect}[1]{\boldsymbol{#1}}
\newcommand{\mat}[1]{\mathbf{#1}}
\newcommand{\diag}{\operatorname{diag}}

\title{Extended Kalman Filter, Vehicle Dynamics, and Odometry Math\\for the RC Car Project}
\author{Generated from the current repository implementation}
\date{\today}

\begin{document}
\maketitle
\tableofcontents
\newpage

\section{Purpose and Big Picture}

This document explains the math currently implemented in \texttt{ExtendedKalmanFilter.py}, \texttt{DynamicsModel.py}, and \texttt{Odometry.py}, using the active values in the \texttt{config/} folder.

At a high level:

\begin{itemize}
    \item the \textbf{dynamics model} predicts how the RC car should move
    \item the \textbf{sensor model} predicts what GPS and IMU should measure from a given state
    \item the \textbf{EKF} combines both while tracking uncertainty
\end{itemize}

The nonlinear filtering equations are:

\begin{align}
\vect{x}_k &= f(\vect{x}_{k-1}, \vect{u}_k) + \vect{w}_k \\
\vect{z}_k &= h(\vect{x}_k) + \vect{v}_k
\end{align}

with process noise and measurement noise:

\begin{align}
\vect{w}_k &\sim \mathcal{N}(\vect{0}, \mat{Q}_k) \\
\vect{v}_k &\sim \mathcal{N}(\vect{0}, \mat{R}_k)
\end{align}

The EKF alternates between:

\begin{itemize}
    \item \textbf{predict}: use the dynamics model to propagate the state
    \item \textbf{update}: use sensors to correct the propagated state
\end{itemize}

\section{Frames, Signs, and Config Values}

From \texttt{config/frames.yaml} and the code:

\begin{itemize}
    \item body frame: $+X$ forward, $+Y$ left, $+Z$ up
    \item world frame: local ENU, so $x = \text{East}$ and $y = \text{North}$
    \item yaw $\psi$ is radians, positive counter-clockwise from East
    \item steering sign convention is currently \texttt{negative-left}
\end{itemize}

That steering convention means:

\begin{equation}
\delta_{\text{math}} = -\delta_{\text{logged}}
\end{equation}

The active vehicle values from \texttt{config/vehicle.yaml} are:

\begin{center}
\begin{tabular}{lll}
\toprule
Symbol & Meaning & Value \\
\midrule
$m$ & mass & $2.667~\text{kg}$ \\
$L$ & wheelbase & $0.3048~\text{m}$ \\
$a$ & front CG distance & $0.159512~\text{m}$ \\
$b$ & rear CG distance & $0.140208~\text{m}$ \\
$t$ & track width & $0.3048~\text{m}$ \\
$h$ & CG height & $0.066309~\text{m}$ \\
$I_z$ & yaw inertia & $0.6~\text{kg}\,\text{m}^2$ \\
$C_f$ & front cornering stiffness & $60.0~\text{N/rad}$ \\
$C_r$ & rear cornering stiffness & $35.0~\text{N/rad}$ \\
$C_d$ & longitudinal drag coefficient & $0.0$ \\
$v_{x,\min}$ & minimum safe longitudinal speed & $0.3~\text{m/s}$ \\
\bottomrule
\end{tabular}
\end{center}

The active noise values from \texttt{config/ekf.yaml} and \texttt{config/sensors.yaml} are:

\begin{center}
\begin{tabular}{lll}
\toprule
Parameter & Meaning & Value \\
\midrule
$\sigma_{\text{gps,pos}}$ & GPS position sigma & $2.5~\text{m}$ \\
$\sigma_{\text{gps,vel}}$ & GPS velocity sigma & $0.05~\text{m/s}$ \\
$\sigma_{\psi}$ & IMU yaw sigma & $0.026~\text{rad}$ \\
$\sigma_r$ & IMU yaw-rate sigma & $0.0014~\text{rad/s}$ \\
$\sigma_{\text{model,vel}}$ & velocity process sigma & $1.5~\text{m/s}^2$ \\
$\sigma_{\text{model,yaw}}$ & yaw-accel process sigma & $2.0~\text{rad/s}^2$ \\
$\sigma_{b_a}$ & accel-bias walk sigma & $0.05~\text{m/s}^3$ \\
$\sigma_{b_g}$ & gyro-bias walk sigma & $0.01~\text{rad/s}^2$ \\
\bottomrule
\end{tabular}
\end{center}

\section{State, Input, and Measurement Definitions}

The EKF state in \texttt{ExtendedKalmanFilter.py} is:

\begin{equation}
\vect{x} =
\begin{bmatrix}
p_e \\ p_n \\ \psi \\ v_x \\ v_y \\ r \\ b_{ax} \\ b_{gz}
\end{bmatrix}
\end{equation}

where:

\begin{itemize}
    \item $p_e, p_n$ are East and North position
    \item $\psi$ is yaw angle
    \item $v_x, v_y$ are body-frame longitudinal and lateral velocity
    \item $r$ is yaw rate
    \item $b_{ax}$ is accelerometer bias
    \item $b_{gz}$ is gyro z-axis bias
\end{itemize}

The predict-step input is:

\begin{equation}
\vect{u} =
\begin{bmatrix}
a_{x,\text{meas}} \\ \delta_{\text{logged}}
\end{bmatrix}
\end{equation}

The active measurement types are:

\begin{itemize}
    \item GPS position
    \item GPS velocity
    \item IMU yaw
    \item IMU yaw rate
\end{itemize}

\section{Dynamics Model}

\subsection{Body-to-world rotation}

The code uses the planar rotation:

\begin{equation}
\mat{R}(\psi) =
\begin{bmatrix}
\cos\psi & -\sin\psi \\
\sin\psi & \cos\psi
\end{bmatrix}
\end{equation}

so the world velocity is:

\begin{equation}
\begin{bmatrix}
v_e \\ v_n
\end{bmatrix}
=
\mat{R}(\psi)
\begin{bmatrix}
v_x \\ v_y
\end{bmatrix}
=
\begin{bmatrix}
v_x\cos\psi - v_y\sin\psi \\
v_x\sin\psi + v_y\cos\psi
\end{bmatrix}
\end{equation}

Therefore:

\begin{align}
\dot{p}_e &= v_x\cos\psi - v_y\sin\psi \\
\dot{p}_n &= v_x\sin\psi + v_y\cos\psi \\
\dot{\psi} &= r
\end{align}

\subsection{Slip angles}

The dynamics model protects against division by zero with:

\begin{equation}
v_{x,\text{safe}} = \operatorname{sign}(v_x)\max(|v_x|, 0.3)
\end{equation}

The slip-angle equations are:

\begin{align}
\alpha_f &= \delta_{\text{math}} - \frac{v_y + ar}{v_{x,\text{safe}}} \\
\alpha_r &= -\frac{v_y - br}{v_{x,\text{safe}}}
\end{align}

Using the active values:

\begin{align}
\alpha_f &= \delta_{\text{math}} - \frac{v_y + 0.159512\,r}{v_{x,\text{safe}}} \\
\alpha_r &= -\frac{v_y - 0.140208\,r}{v_{x,\text{safe}}}
\end{align}

\subsection{Lateral tire forces}

The code uses a linear tire model:

\begin{align}
F_{yf} &= 2C_f\alpha_f = 120\,\alpha_f \\
F_{yr} &= 2C_r\alpha_r = 70\,\alpha_r
\end{align}

\subsection{Bias-corrected acceleration}

The IMU longitudinal acceleration is corrected by the EKF bias estimate:

\begin{equation}
a_{x,\text{corr}} = a_{x,\text{meas}} - b_{ax}
\end{equation}

\subsection{Continuous-time nonlinear state derivative}

The model implemented in \texttt{DynamicsModel.py} is:

\begin{align}
\dot{p}_e &= v_x\cos\psi - v_y\sin\psi \\
\dot{p}_n &= v_x\sin\psi + v_y\cos\psi \\
\dot{\psi} &= r \\
\dot{v}_x &= a_{x,\text{corr}} + v_y r - C_d v_x|v_x| \\
\dot{v}_y &= -v_x r + \frac{F_{yf}\cos\delta_{\text{math}} + F_{yr}}{m} \\
\dot{r} &= \frac{aF_{yf}\cos\delta_{\text{math}} - bF_{yr}}{I_z} \\
\dot{b}_{ax} &= 0 \\
\dot{b}_{gz} &= 0
\end{align}

Because $C_d = 0$ in the current config, the longitudinal equation simplifies to:

\begin{equation}
\dot{v}_x = a_{x,\text{corr}} + v_y r
\end{equation}

Substituting config values into the lateral and yaw equations:

\begin{align}
\dot{v}_y &= -v_x r + \frac{120\,\alpha_f\cos\delta_{\text{math}} + 70\,\alpha_r}{2.667} \\
\dot{r} &= \frac{0.159512\cdot 120\,\alpha_f\cos\delta_{\text{math}} - 0.140208\cdot 70\,\alpha_r}{0.6}
\end{align}

\subsection{Discrete-time process model}

The EKF uses forward Euler discretization:

\begin{equation}
\vect{x}_k^- = f_d(\vect{x}_{k-1}^+, \vect{u}_k)
\approx \vect{x}_{k-1}^+ + f(\vect{x}_{k-1}^+, \vect{u}_k)\Delta t
\end{equation}

and then wraps yaw:

\begin{equation}
\psi \leftarrow \operatorname{atan2}(\sin\psi,\cos\psi)
\end{equation}

\section{What the F Matrix Is}

The matrix $\mat{F}_k$ is the \textbf{state-transition Jacobian}. It is the local linearization of the discrete nonlinear process model:

\begin{equation}
\mat{F}_k = \left.\frac{\partial f_d}{\partial \vect{x}}\right|_{\vect{x}_{k-1}^+,\,\vect{u}_k}
\end{equation}

In plain language:

\begin{quote}
If the current state estimate changes a little, how does the next predicted state change?
\end{quote}

That is why $\mat{F}_k$ appears inside the covariance prediction:

\begin{equation}
\mat{P}_k^- = \mat{F}_k \mat{P}_{k-1}^+ \mat{F}_k^T + \mat{Q}_k
\end{equation}

The repository computes $\mat{F}_k$ numerically with a central-difference Jacobian:

\begin{equation}
\mat{F}_k[:,i] =
\frac{
f_d(\vect{x} + \varepsilon\vect{e}_i,\vect{u}) -
f_d(\vect{x} - \varepsilon\vect{e}_i,\vect{u})
}{2\varepsilon}
\end{equation}

with $\varepsilon = 10^{-6}$.

\subsection{Analytic continuous-time Jacobian for intuition}

Ignoring yaw wrapping and the low-speed clamp for a moment, and assuming $v_x > v_{x,\min}$, the continuous-time Jacobian $\mat{A} = \partial f / \partial \vect{x}$ has the sparse structure:

\begin{equation}
\mat{A} =
\begin{bmatrix}
0 & 0 & a_{13} & \cos\psi & -\sin\psi & 0 & 0 & 0 \\
0 & 0 & a_{23} & \sin\psi & \cos\psi & 0 & 0 & 0 \\
0 & 0 & 0 & 0 & 0 & 1 & 0 & 0 \\
0 & 0 & 0 & a_{44} & r & v_y & -1 & 0 \\
0 & 0 & 0 & a_{54} & a_{55} & a_{56} & 0 & 0 \\
0 & 0 & 0 & a_{64} & a_{65} & a_{66} & 0 & 0 \\
0 & 0 & 0 & 0 & 0 & 0 & 0 & 0 \\
0 & 0 & 0 & 0 & 0 & 0 & 0 & 0
\end{bmatrix}
\end{equation}

where:

\begin{align}
a_{13} &= -v_x\sin\psi - v_y\cos\psi \\
a_{23} &= v_x\cos\psi - v_y\sin\psi \\
a_{44} &= -2C_d|v_x| \\
a_{54} &= -r + \frac{2C_f\cos\delta_{\text{math}}(v_y + ar) + 2C_r(v_y - br)}{m v_x^2} \\
a_{55} &= \frac{-2C_f\cos\delta_{\text{math}} - 2C_r}{m v_x} \\
a_{56} &= -v_x + \frac{-2aC_f\cos\delta_{\text{math}} + 2bC_r}{m v_x} \\
a_{64} &= \frac{2aC_f\cos\delta_{\text{math}}(v_y + ar) - 2bC_r(v_y - br)}{I_z v_x^2} \\
a_{65} &= \frac{-2aC_f\cos\delta_{\text{math}} + 2bC_r}{I_z v_x} \\
a_{66} &= \frac{-2a^2C_f\cos\delta_{\text{math}} - 2b^2C_r}{I_z v_x}
\end{align}

For forward Euler discretization:

\begin{equation}
\mat{F}_k \approx \mat{I} + \mat{A}_k \Delta t
\end{equation}

The code does not use this symbolic matrix directly, but it is a useful way to understand what the numerical Jacobian is approximating.

\section{Predict Step}

The EKF predict step in \texttt{ExtendedKalmanFilter.py} is:

\begin{align}
\vect{x}_k^- &= f_d(\vect{x}_{k-1}^+, \vect{u}_k) \\
\mat{F}_k &= \left.\frac{\partial f_d}{\partial \vect{x}}\right|_{\vect{x}_{k-1}^+,\,\vect{u}_k} \\
\mat{P}_k^- &= \mat{F}_k \mat{P}_{k-1}^+ \mat{F}_k^T + \mat{Q}_k
\end{align}

Here:

\begin{itemize}
    \item $\vect{x}_k^-$ is the predicted state
    \item $\mat{P}_k^-$ is the predicted covariance
    \item $\mat{Q}_k$ is the process-noise covariance
\end{itemize}

\subsection{Initial covariance}

From \texttt{config/ekf.yaml}, the initial covariance is:

\begin{equation}
\mat{P}_0 = \diag(
10.0^2,\,
10.0^2,\,
0.5235987756^2,\,
2.0^2,\,
2.0^2,\,
0.3490658504^2,\,
1.0^2,\,
0.0872664626^2)
\end{equation}

Numerically:

\begin{equation}
\mat{P}_0 = \diag(
100.0,\,
100.0,\,
0.2741556778,\,
4.0,\,
4.0,\,
0.1218469679,\,
1.0,\,
0.0076154355)
\end{equation}

\subsection{Process-noise matrix}

The code constructs:

\begin{align}
\sigma_v^2 &= (1.5\Delta t)^2 \\
\sigma_{\psi}^2 &= (2.0\Delta t)^2 \\
\sigma_{b_a}^2 &= (0.05\Delta t)^2 \\
\sigma_{b_g}^2 &= (0.01\Delta t)^2
\end{align}

so that:

\begin{equation}
\mat{Q}(\Delta t) =
\diag(
10^{-4},\,
10^{-4},\,
(2.0\Delta t)^2,\,
(1.5\Delta t)^2,\,
(1.5\Delta t)^2,\,
(2.0\Delta t)^2,\,
(0.05\Delta t)^2,\,
(0.01\Delta t)^2)
\end{equation}

For $\Delta t = 0.1$ seconds:

\begin{equation}
\mat{Q}(0.1) =
\diag(
10^{-4},\,
10^{-4},\,
4.0\times 10^{-2},\,
2.25\times 10^{-2},\,
2.25\times 10^{-2},\,
4.0\times 10^{-2},\,
2.5\times 10^{-5},\,
10^{-6})
\end{equation}

\section{Sensor Models and H Matrices}

Each measurement channel defines:

\begin{equation}
\vect{z}_k = h(\vect{x}_k) + \vect{v}_k
\end{equation}

and a measurement Jacobian:

\begin{equation}
\mat{H}_k = \left.\frac{\partial h}{\partial \vect{x}}\right|_{\vect{x}_k^-}
\end{equation}

\subsection{GPS position}

The GPS position model is:

\begin{equation}
h_{\text{gps,pos}}(\vect{x}) =
\begin{bmatrix}
p_e \\ p_n
\end{bmatrix}
\end{equation}

with:

\begin{equation}
\mat{H}_{\text{gps,pos}} =
\begin{bmatrix}
1 & 0 & 0 & 0 & 0 & 0 & 0 & 0 \\
0 & 1 & 0 & 0 & 0 & 0 & 0 & 0
\end{bmatrix}
\end{equation}

and default covariance:

\begin{equation}
\mat{R}_{\text{gps,pos}} = \diag(2.5^2, 2.5^2) = \diag(6.25, 6.25)
\end{equation}

If GPS reports $e_{px}$ and $e_{py}$, the code uses:

\begin{equation}
\mat{R}_{\text{gps,pos}} = \diag(e_{px}^2, e_{py}^2)
\end{equation}

\subsection{GPS velocity}

Because the state stores body-frame velocity but GPS reports world-frame velocity:

\begin{equation}
h_{\text{gps,vel}}(\vect{x}) =
\begin{bmatrix}
v_x\cos\psi - v_y\sin\psi \\
v_x\sin\psi + v_y\cos\psi
\end{bmatrix}
\end{equation}

The code computes the Jacobian numerically, but the analytic form is:

\begin{equation}
\mat{H}_{\text{gps,vel}} =
\begin{bmatrix}
0 & 0 & -v_x\sin\psi - v_y\cos\psi & \cos\psi & -\sin\psi & 0 & 0 & 0 \\
0 & 0 & v_x\cos\psi - v_y\sin\psi & \sin\psi & \cos\psi & 0 & 0 & 0
\end{bmatrix}
\end{equation}

The default covariance is:

\begin{equation}
\mat{R}_{\text{gps,vel}} = \diag(0.05^2, 0.05^2) = \diag(0.0025, 0.0025)
\end{equation}

If GPS reports $e_{ps}$, the code uses:

\begin{equation}
\mat{R}_{\text{gps,vel}} = \diag(e_{ps}^2, e_{ps}^2)
\end{equation}

This update is skipped when the reported speed is below the configured threshold:

\begin{equation}
\text{speed} < 0.2~\text{m/s}
\end{equation}

\subsection{IMU yaw}

The yaw model is:

\begin{equation}
h_{\psi}(\vect{x}) =
\begin{bmatrix}
\psi
\end{bmatrix}
\end{equation}

with:

\begin{equation}
\mat{H}_{\psi} =
\begin{bmatrix}
0 & 0 & 1 & 0 & 0 & 0 & 0 & 0
\end{bmatrix}
\end{equation}

and covariance:

\begin{equation}
\mat{R}_{\psi} =
\begin{bmatrix}
0.026^2
\end{bmatrix}
=
\begin{bmatrix}
0.000676
\end{bmatrix}
\end{equation}

The innovation is wrapped to keep angle errors in $[-\pi,\pi)$.

\subsection{IMU yaw rate}

The gyro z-axis measurement model is:

\begin{equation}
h_r(\vect{x}) =
\begin{bmatrix}
r + b_{gz}
\end{bmatrix}
\end{equation}

with:

\begin{equation}
\mat{H}_{r} =
\begin{bmatrix}
0 & 0 & 0 & 0 & 0 & 1 & 0 & 1
\end{bmatrix}
\end{equation}

and covariance:

\begin{equation}
\mat{R}_{r} =
\begin{bmatrix}
0.0014^2
\end{bmatrix}
=
\begin{bmatrix}
1.96\times 10^{-6}
\end{bmatrix}
\end{equation}

\section{Update Step}

For any measurement update, the EKF computes:

\begin{align}
\vect{y}_k &= \vect{z}_k - h(\vect{x}_k^-) \\
\mat{S}_k &= \mat{H}_k \mat{P}_k^- \mat{H}_k^T + \mat{R}_k \\
\mat{K}_k &= \mat{P}_k^- \mat{H}_k^T \mat{S}_k^{-1} \\
\vect{x}_k^+ &= \vect{x}_k^- + \mat{K}_k \vect{y}_k
\end{align}

where:

\begin{itemize}
    \item $\vect{y}_k$ is the innovation or residual
    \item $\mat{S}_k$ is the innovation covariance
    \item $\mat{K}_k$ is the Kalman gain
\end{itemize}

The code uses the Joseph-form covariance update:

\begin{equation}
\mat{P}_k^+ =
(\mat{I} - \mat{K}_k\mat{H}_k)\mat{P}_k^-(\mat{I} - \mat{K}_k\mat{H}_k)^T
+ \mat{K}_k\mat{R}_k\mat{K}_k^T
\end{equation}

This is more numerically stable than the minimal form $ \mat{P}^+ = (\mat{I} - \mat{K}\mat{H})\mat{P}^- $.

\section{GPS Initialization and Coordinate Conversion}

The filter defines a local ENU origin from the first good GPS fix and converts latitude and longitude to local meters with:

\begin{align}
\text{meters\_per\_degree}_{\text{lat}} &= 111320 \\
\text{meters\_per\_degree}_{\text{lon}} &= 111320\cos(\text{lat}_0)
\end{align}

\begin{align}
p_e &= (\text{lon} - \text{lon}_0)\,\text{meters\_per\_degree}_{\text{lon}} \\
p_n &= (\text{lat} - \text{lat}_0)\,\text{meters\_per\_degree}_{\text{lat}}
\end{align}

GPS track is converted into ENU yaw with:

\begin{equation}
\psi = \operatorname{wrap}\left(\frac{\pi}{2} - \operatorname{deg2rad}(\text{track}_{\deg})\right)
\end{equation}

Then world velocity is:

\begin{equation}
\begin{bmatrix}
v_e \\ v_n
\end{bmatrix}
=
\begin{bmatrix}
\text{speed}\cos\psi \\
\text{speed}\sin\psi
\end{bmatrix}
\end{equation}

If needed, the initial body velocity is obtained from the inverse rotation:

\begin{equation}
\begin{bmatrix}
v_x \\ v_y
\end{bmatrix}
=
\begin{bmatrix}
\cos\psi & \sin\psi \\
-\sin\psi & \cos\psi
\end{bmatrix}
\begin{bmatrix}
v_e \\ v_n
\end{bmatrix}
\end{equation}

\section{Load Transfer and Handling}

The dynamics file also computes first-pass body accelerations and load transfer:

\begin{align}
A_x &= a_{x,\text{meas}} - b_{ax} \\
A_y &= \frac{F_{yf}\cos\delta_{\text{math}} + F_{yr}}{m}
\end{align}

\begin{align}
W &= mg \\
\Delta W_x &= W\left(\frac{A_x h}{L}\right) \\
\Delta W_y &= W\left(\frac{A_y h}{t}\right)
\end{align}

Using the active mass:

\begin{equation}
W = 2.667 \cdot 9.80665 = 26.1543~\text{N}
\end{equation}

Since $L = t = 0.3048~\text{m}$, the current scaling becomes:

\begin{align}
\Delta W_x &\approx 5.687\,A_x \\
\Delta W_y &\approx 5.687\,A_y
\end{align}

The handling heuristic is:

\begin{itemize}
    \item understeer if $|\alpha_f| > |\alpha_r| + 1^\circ$
    \item oversteer if $|\alpha_r| > |\alpha_f| + 1^\circ$
    \item otherwise neutral
\end{itemize}

The first-pass Ackermann steering target is:

\begin{equation}
\delta_{\text{ack}} = \arctan\left(\frac{L}{R}\right) = \arctan\left(\frac{0.3048}{R}\right)
\end{equation}

\section{Odometry}

\texttt{Odometry.py} is not a Kalman filter. It is direct IMU-based dead reckoning without a covariance matrix.

The odometry flow is:

\begin{enumerate}
    \item read IMU gyro, acceleration, quaternion, gravity, and Euler angles
    \item optionally smooth linear acceleration
    \item rotate body-frame acceleration into the world frame
    \item integrate acceleration to velocity
    \item integrate velocity to position
\end{enumerate}

If the quaternion rotation matrix is $\mat{R}_{bw}$, then:

\begin{equation}
\vect{a}_{\text{world}} = \mat{R}_{bw}\vect{a}_{\text{body}}
\end{equation}

Velocity is updated with the trapezoidal rule:

\begin{equation}
\vect{v}_k = \vect{v}_{k-1} + \frac{1}{2}\left(\vect{a}_k + \vect{a}_{k-1}\right)\Delta t
\end{equation}

Position is updated the same way:

\begin{equation}
\vect{p}_k = \vect{p}_{k-1} + \frac{1}{2}\left(\vect{v}_k + \vect{v}_{k-1}\right)\Delta t
\end{equation}

The north estimate helper normalizes gravity and magnetometer vectors:

\begin{align}
\hat{\vect{g}} &= \frac{\vect{g}}{\|\vect{g}\|} \\
\hat{\vect{m}} &= \frac{\vect{m}}{\|\vect{m}\|}
\end{align}

then forms:

\begin{align}
\vect{e}_{\text{east}} &= \hat{\vect{g}} \times \hat{\vect{m}} \\
\vect{e}_{\text{north}} &= \vect{e}_{\text{east}} \times \hat{\vect{g}}
\end{align}

\section{Pseudocode}

\subsection{Startup}

\begin{lstlisting}
load config bundle
build DynamicsModel from config
build ExtendedKalmanFilter from config
wait for first good GPS fix
set local ENU origin
initialize state x0 and covariance P0
\end{lstlisting}

\subsection{Main EKF loop}

\begin{lstlisting}
repeat:
    dt = t_now - t_prev
    read steering, IMU accel x, IMU gyro z, optional yaw, GPS
    u_k = [a_x_meas, delta_logged]^T

    predict:
        x^- = f_d(x^+, u_k)
        F_k = d f_d / d x
        P^- = F_k P^+ F_k^T + Q_k

    if GPS position is available:
        update GPS position

    if GPS speed and track are available and speed >= threshold:
        update GPS velocity

    if yaw is available:
        update yaw

    update yaw rate from gyro z
    publish x^+ and P^+
\end{lstlisting}

\subsection{Generic update}

\begin{lstlisting}
y = z - h(x^-)
if angle residual:
    y = wrap(y)

S = H P^- H^T + R
K = P^- H^T S^-1
x^+ = x^- + K y
P^+ = (I - K H) P^- (I - K H)^T + K R K^T
\end{lstlisting}

\section{Most Important Matrices at a Glance}

\begin{align}
\vect{x} &= [p_e,\ p_n,\ \psi,\ v_x,\ v_y,\ r,\ b_{ax},\ b_{gz}]^T \\
\vect{u} &= [a_{x,\text{meas}},\ \delta_{\text{logged}}]^T
\end{align}

\begin{equation}
\mat{R}(\psi) =
\begin{bmatrix}
\cos\psi & -\sin\psi \\
\sin\psi & \cos\psi
\end{bmatrix}
\end{equation}

\begin{equation}
\mat{P}_0 = \diag(100,\ 100,\ 0.2741556778,\ 4,\ 4,\ 0.1218469679,\ 1,\ 0.0076154355)
\end{equation}

\begin{equation}
\mat{Q}(\Delta t) = \diag(10^{-4},\ 10^{-4},\ (2\Delta t)^2,\ (1.5\Delta t)^2,\ (1.5\Delta t)^2,\ (2\Delta t)^2,\ (0.05\Delta t)^2,\ (0.01\Delta t)^2)
\end{equation}

\begin{equation}
\mat{H}_{\text{gps,pos}} =
\begin{bmatrix}
1 & 0 & 0 & 0 & 0 & 0 & 0 & 0 \\
0 & 1 & 0 & 0 & 0 & 0 & 0 & 0
\end{bmatrix}
\end{equation}

\begin{equation}
\mat{H}_{\psi} =
\begin{bmatrix}
0 & 0 & 1 & 0 & 0 & 0 & 0 & 0
\end{bmatrix}
\end{equation}

\begin{equation}
\mat{H}_{r} =
\begin{bmatrix}
0 & 0 & 0 & 0 & 0 & 1 & 0 & 1
\end{bmatrix}
\end{equation}

\begin{equation}
\mat{R}_{\text{gps,pos}} = \diag(6.25, 6.25)
\end{equation}

\begin{equation}
\mat{R}_{\text{gps,vel}} = \diag(0.0025, 0.0025)
\end{equation}

\begin{equation}
\mat{R}_{\psi} =
\begin{bmatrix}
0.000676
\end{bmatrix}
\end{equation}

\begin{equation}
\mat{R}_{r} =
\begin{bmatrix}
1.96\times 10^{-6}
\end{bmatrix}
\end{equation}

\section{Building the PDF}

Use the helper script from the repository root:

\begin{lstlisting}
powershell -ExecutionPolicy Bypass -File .\build_ekf_pdf.ps1
\end{lstlisting}

It will:

\begin{itemize}
    \item use \texttt{pdflatex} if installed
    \item otherwise use \texttt{tectonic} if installed
    \item otherwise stop with a clear error
\end{itemize}

In the current environment, no local LaTeX engine was installed, so the PDF could not be built automatically here.

\section{Conclusion}

The current repository structure is:

\begin{itemize}
    \item \texttt{DynamicsModel.py}: provides the nonlinear process model $f(\vect{x},\vect{u})$
    \item \texttt{ExtendedKalmanFilter.py}: wraps that model in EKF predict and update steps
    \item \texttt{Odometry.py}: provides direct IMU-based dead reckoning without covariance
\end{itemize}

If someone asks, ``What is the $F$ matrix again?'', the most useful answer is:

\begin{quote}
$F_k$ is the matrix that tells the EKF how uncertainty in the current state propagates through the nonlinear predict step into uncertainty in the next predicted state.
\end{quote}

\end{document}
